\section{The Mechanized Proof}
\label{app:lean}

This appendix gives the formal definitions, the decision chain, and the
correspondence with the paper. Names are in namespace \texttt{DeciNSSE}
unless qualified otherwise. Listings suppress proof bodies and implicit
variable declarations where these do not affect the statement. Assignments
in the source may be named $\rho$, $\sigma$, or $A$; the paper uses variants
of $\sigma$ for assignments. Paths are lists over
\texttt{Fin n}; binary examples use $\{0,1\}$.

\subsection{Trusted Definitions and Their Characterization}
\label{app:lean-statement}
The primitive symbols are \texttt{Sym.bot}, \texttt{Sym.f}, and
\texttt{Sym.top}, with their three-element order, induced by \lean{Sym.rank} with values
$0,1,2$. The trusted statement also includes \lean{polarity}, the parity
of contravariant positions on a path, and the tree operations
\lean{Tree.node}, \lean{Tree.bot}, and \lean{Tree.top}.
Trees and their variance order are defined by
\begin{lstlisting}
structure Tree (n : ℕ) where
  fn : List (Fin n) → Option Sym
  wf : fn [] ≠ none ∧
       ∀ π i, (fn (π ++ [i])).isSome ↔ fn π = some Sym.f

def Tree.Le (c : Fin n → Bool) (t u : Tree n) : Prop :=
  ∀ π a b, t.fn π = some a → u.fn π = some b →
    if polarity c π then b ≤ a else a ≤ b
\end{lstlisting}
Here \texttt{true} is contravariant. The notation $\le$ on \texttt{Tree n}
means the internal covariant order \lean{Tree.CovLe}; the public statements
use \lean{Tree.Le} explicitly.

The module \texttt{Semantics/Characterization} relates these definitions
to conventional semantics. \lean{Characterization.treeLe_eq_gfp} identifies
the order with the greatest fixed point of the textbook rule;
\lean{Characterization.lfp_ne_treeLe} distinguishes the least fixed point
at positive arity. \lean{Characterization.unfold_range_iff} characterizes
the range of graph unfolding by finiteness of the set of rooted subtrees,
and \lean{Characterization.toTree_range_iff} characterizes the image of
finite terms by finite tree domain. The interpretation maps \lean{RGraph.unfold} and \lean{FTree.toTree}
are themselves part of the trusted statement; the two range theorems
justify quantification over these representations of the intended tree classes. The module also proves the precise
flat expressiveness distinction:
\lean{Characterization.le_definable_iff} and
\lean{Characterization.nullary_not_definable} cover positive and zero arity.
General flattening of nested input syntax is not part of the exported API.

Variables are \lean{V}, an abbreviation for \texttt{Fin k}, and
\lean{Constraint} is a list of the following literals:
\begin{lstlisting}
inductive Lit (n k : ℕ) where
  | leF (x : V k) (children : Fin n → V k)
  | fLe (children : Fin n → V k) (x : V k)
  | eqBot (x : V k)
  | eqTop (x : V k)

def Sat (c : Fin n → Bool) (ρ : V k → Tree n)
    (ϕ : Constraint n k) : Prop :=
  ∀ l ∈ ϕ, Lit.holds c ρ l

abbrev Satisfiable (c : Fin n → Bool) (ϕ : Constraint n k) : Prop :=
  ∃ ρ, Sat c ρ ϕ

def Entails (c : Fin n → Bool) (ϕ : Constraint n k) (x y : V k) : Prop :=
  ∀ ρ, Sat c ρ ϕ → Tree.Le c (ρ x) (ρ y)
\end{lstlisting}
The interpretation \lean{Lit.holds} uses constructor trees, the explicit
variance order, and equality for the constant literals. Satisfaction by
an assignment is distinct from existence of a solution.
\begin{lstlisting}
def EntailsReg (c : Fin n → Bool) (ϕ : Constraint n k) (x y : V k) : Prop :=
  ∀ σ : V k → RGraph n, Sat c (RGraph.unfold ∘ σ) ϕ →
    Tree.Le c (σ x).unfold (σ y).unfold

def EntailsFin (c : Fin n → Bool) (ϕ : Constraint n k) (x y : V k) : Prop :=
  ∀ σ : V k → FTree n, Sat c (FTree.toTree ∘ σ) ϕ →
    Tree.Le c (σ x).toTree (σ y).toTree
\end{lstlisting}

\subsection{Public API and Module Layout}
\label{app:lean-artifact}
Importing \texttt{DeciNSSE} exposes the three procedures:
\begin{lstlisting}
def decideEntails {n k : ℕ} (c : Fin n → Bool)
    (ϕ : Constraint n k) (x y : V k) : Decidable (Entails c ϕ x y)
def decideEntailsReg {n k : ℕ} (c : Fin n → Bool)
    (ϕ : Constraint n k) (x y : V k) : Decidable (EntailsReg c ϕ x y)
def decideEntailsFin {n k : ℕ} (c : Fin n → Bool)
    (ϕ : Constraint n k) (x y : V k) : Decidable (EntailsFin c ϕ x y)

theorem decideEntails_correct {n k : ℕ} (c : Fin n → Bool)
    (ϕ : Constraint n k) (x y : V k) :
  @decide _ (decideEntails c ϕ x y) = true ↔ Entails c ϕ x y
\end{lstlisting}
The other correctness declarations are \lean{decideEntailsReg_correct} and
\lean{decideEntailsFin_correct}. None assumes satisfiability, positive
arity, an inhabited path alphabet, or a supplied reader family.
\lean{productVariance} and \lean{arrowVariance}, in \texttt{Instances},
provide the two named binary special cases. Importing \texttt{DeciNSSE.Main}
suffices for the three procedures. The summary and unary procedures of
\cref{app:lean-complexity} have the same types, the unary ones at arity one.
All of them instantiate
\lean{Monitor.decideEntailsOfMonitors} and its regular and finite versions,
which take the two side monitors and a decision procedure for side holes.

The artifact also contains a certifying solver, \lean{solve}. It searches for
an entailment certificate, which records unsatisfiability or a closed cover
of each query side, and for a refutation, which records a side and a hole;
it checks every candidate with proved checkers and falls back to
\lean{decideEntails}. \lean{solve_correct} states
that its answer is \texttt{true} iff \lean{Entails} holds, and a refutation
carries a regular countermodel. Under \texttt{import DeciNSSE}, compiled code
of \lean{decideEntails} runs the packed decider
\lean{Solver.decideEntailsPacked}, by the \texttt{csimp} lemma
\lean{Solver.decideEntails_eq_decideEntailsPacked}; no other procedure is
redirected, and the kernel-checked statements are unaffected.

\begin{center}\small
\begin{tabular}{@{}lp{0.70\linewidth}@{}}
\toprule
Module group & Contents\\\midrule
\texttt{Words} & Finite-alphabet paths and periodicity.\\
\texttt{Semantics} & Trees, finite terms, regular graphs, variance, normalization, polarity selection, and characterization.\\
\texttt{Constraints} & Flat syntax, satisfaction, signed translation, closure, and path bounds.\\
\texttt{Satisfiability} & Least shapes, graph realization, clash tests, finite and regular solutions.\\
\texttt{Transfer} & Query sides and unsafe words, ranked extensions, cycle reflection, and countermodel transfers.\\
\texttt{Monitor} & Four-spine extension, closure and clashes, events, label reader, side monitors, semantic interface, and the generic procedures.\\
\texttt{Holes} & Readers and holes; the canonical hierarchy (\texttt{Holes/Hierarchy/*}), depth, cores, and supports; desubstitution, co-alignment, exchange types, level maps, normal forms, and boundary pairs; agreement, the column automaton, bounded-depth compression, and the elementary bound.\\
\texttt{RejectedTail} & Rejected-path property, branching and forced tails, return gap, depth, and finite search.\\
\texttt{Instances}, \texttt{Main} & Named variance maps and the three public decision procedures.\\
\texttt{Countermodels} & The Boolean refutation check and computable regular and finite countermodels.\\
\texttt{Certificates} & Certificate formats, covering products, the specification checkers and their soundness, and the JSON front end of the certificate checker.\\
\texttt{Solver} & The certifying solver: packed readers and checks, the packed decider, certificate search, reachable countermodels, and the JSON front end of the solver.\\
\texttt{Complexity} & The summary search, its configuration bound, and its procedures; one-letter readers and the unary procedures; binary automata and shared tree constructions; the hardness reduction; the counter automaton and the systems with exponentially long holes.\\\bottomrule
\end{tabular}
\end{center}
The executable front ends \texttt{Certificates/Command} and
\texttt{Solver/Command} and their JSON parser \texttt{Certificates/Json} are
built by the executable targets and are not reachable from
\texttt{import DeciNSSE}; every other module is.

The files \texttt{lean-toolchain} and \texttt{lakefile.toml} pin Lean and
Mathlib, respectively, to \texttt{v4.35.0}. To check the artifact, run
\begin{lstlisting}
lake exe cache get
lake build DeciNSSE
\end{lstlisting}
The first command downloads compiled dependencies. The kernel checks the
proofs, including the reduction to finite search; compiled execution also
relies on the compiler. The three procedures are computable definitions,
not classical choices of decidability instances. Proofs may use classical
reasoning. The package's axiom audit reports only \texttt{propext},
\texttt{Classical.choice}, and \texttt{Quot.sound} for the procedures,
their correctness theorems, the solver's correctness theorem, and the
headline theorems of \texttt{Complexity}. These dependencies can be inspected with
\begin{lstlisting}
import DeciNSSE
#print axioms DeciNSSE.decideEntails
#print axioms DeciNSSE.decideEntailsReg
#print axioms DeciNSSE.decideEntailsFin
#print axioms DeciNSSE.decideEntails_correct
#print axioms DeciNSSE.decideEntailsReg_correct
#print axioms DeciNSSE.decideEntailsFin_correct
#print axioms DeciNSSE.solve_correct
#print axioms DeciNSSE.decideEntailsSummary_correct
#print axioms DeciNSSE.Complexity.Summary.card_config_label_le
#print axioms DeciNSSE.decideEntailsUnary_correct
#print axioms DeciNSSE.Complexity.Unary.exists_hole_length_le
#print axioms DeciNSSE.Complexity.entailsFin_hardnessSystem_iff
#print axioms DeciNSSE.Complexity.expSystem_shortest_hole
\end{lstlisting}

\subsection{The Decision Chain}
\label{app:lean-decision}
The function \lean{signed} constructs the signed constraint system.
The function \lean{Signed.dual} implements the assignment involution and \lean{SignSelection.select} implements polarity selection on
assignments with a common shape. The shared theorem
\lean{SignSelection.select_leastShape_sat_fixed} states satisfaction and
fixedness of the selected least shape of a flip-closed, sign-coherent
system without label clash. It is used by the witness, regular, and finite
constructions. \lean{Monitor.satB} decides satisfiability by the signed
label-clash test. \lean{Spine.extension} enforces four prefix conditions
in a signed layout directly; \texttt{Monitor/Closure} proves its closure
invariants, and \texttt{Monitor/Clash} and \texttt{Monitor/Semantics}
identify clashes with inherited clashes or events and then with unsafety.
\lean{Events.sideUnsafe_iff_not_labelClash} has no satisfiability premise.

The label reader \lean{Monitor.reader} uses
\begin{lstlisting}
structure Monitor.State (m : ℕ) where
  U : Finset (V m)
  L : Finset (V m)
  b : Bool
\end{lstlisting}
The reader, target, and admission definitions have no side parameter.
\lean{Events.sideSystem} chooses the signed system or its order dual, and
\lean{Events.sideQuery} chooses $(x^+,y^+)$ or $(y^+,x^+)$;
\lean{Monitor.sideReader} and \lean{Monitor.sideTarget} are the reader and
target of a side.
The acceptance latch records readiness at a strictly earlier cut.
\lean{Monitor.target} is the complement of acceptance, and
\lean{Monitor.admission} is the disjunction of cross and self admission.
\lean{Monitor.hole_iff_not_occurs} requires no satisfiability hypothesis;
\lean{Monitor.hole_iff_unsafe} requires a solution, precisely to exclude
inherited clashes. \lean{Monitor.rejectedPath} supplies the rejected-path
property without a satisfiability assumption. \lean{Monitor.family}
packages the two side readers for the given variance.

\lean{RejectedTail.RejectedPath} expresses betweenness and the
admission cone. The general procedure
\lean{RejectedTail.decideHole} enumerates finite-alphabet words up to
\lean{RejectedTail.holeLengthBound}. Its bound is defined by
\begin{lstlisting}
def holeDepthBound (N : ℕ) : ℕ :=
  2 * N * (4*N - 1) - 1

def holeLengthBound (N : ℕ) : ℕ :=
  ElementaryBound.lengthBound N (holeDepthBound N)
\end{lstlisting}
Here $N$ is the original reader-state count; \cref{sec:mechanization}
explains how derivative states enter the bound. The theorem
\lean{RejectedTail.return_gap} holds for every reader and gives
$e-s\le |w|-J$; it assumes neither that $w$ is a hole nor the rejected-path property. Intermediate depth and compression theorems assume an inhabited
alphabet, which nonempty words supply; the final bounded-length theorem
\lean{RejectedTail.hole_iff_bounded_length} removes this assumption and
explicitly covers the empty alphabet.

The bounded-depth theorem uses the predicate \lean{LetteredHierarchy.InL}:
$\mathsf{InL}(d,w)$ means that some canonical level $i\le d$ is unary.
The bound in \eqref{eq:holebound} is \lean{ElementaryBound.lengthBound}:
\begin{lstlisting}
def dictionaryBound (N d : ℕ) : ℕ :=
  2 + NormalForm.normalAllowance N d * 3 ^ d

def lengthBound (N d : ℕ) : ℕ :=
  (N + d + 1) *
    BoundedDepth.descentFactor N d (dictionaryBound N d) ^ d
\end{lstlisting}
The column count is \lean{BoundedDepth.colBound}, and
\lean{BoundedDepth.descentFactor N d g} is that count at $N+d+1$ states
and $g$ tracks. \lean{NormalForm.exists_normal_form} preserves canonical
depth and all start counts, never increases length or a dictionary, and
recomputes the protected exchange types after each merge.
\lean{ElementaryBound.exists_bounded_dictionary_hole} bounds the
positive dictionaries of the resulting word by
\lean{ElementaryBound.dictionaryBound}.

\lean{BoundedDepth.shorten_level_supported} keeps its recursive support
inside the original word's dictionary and deduplicates it. Its identity
base case needs no bound on the level-zero alphabet.
\lean{BoundedDepth.exists_short_hole_of_dictionary} combines this descent
with unary shortening. Finally, \lean{ElementaryBound.exists_short_hole}
composes normalization and descent; it retains membership in
$\mathsf{InL}(d,\cdot)$, not necessarily equality of the first unary level.
The back end invokes this theorem directly at
\lean{RejectedTail.holeDepthBound}, with no intermediate stopping reader.

The arithmetic estimates \lean{ElementaryBound.dictionaryBound_le} and
\lean{ElementaryBound.lengthBound_le} yield the final checked inequality
\begin{lstlisting}
theorem holeLengthBound_le (N : ℕ) :
  holeLengthBound N ≤ 2 ^ 2 ^ (1280 * (N + 1) ^ 5)
\end{lstlisting}
in namespace \texttt{RejectedTail}. The statements hold for every natural
number $N$; actual monitors have a start state and hence $N\ge1$.
The cost of enumerating all words up to this length
(\cref{sec:complexity-upper}) is not mechanized;
\cref{app:lean-complexity} describes the mechanized content of
\cref{sec:complexity}.

The entailment procedure \lean{Monitor.decideEntailsOf} combines
satisfiability testing with the two reader searches, here the enumeration
\lean{Monitor.SideMonitor.decideHoles} up to the length bound. \lean{entailsReg_iff_entails} and \lean{entailsFin_iff}
provide the transfers. The finite procedure first tests \lean{satFinB};
its equivalence with finite satisfiability uses the selector construction
on the signed least-shape solution. \texttt{Main} instantiates the family,
so there are no undischarged interface premises in the public API.

\subsection{Correspondence}
\label{app:lean-map}
\begin{longtable}{@{}>{\raggedright\arraybackslash}p{0.31\linewidth}>{\raggedright\arraybackslash}p{0.65\linewidth}@{}}
\caption{Paper results and declarations in namespace \texttt{DeciNSSE}.
A dagger marks differences in intermediate statement strength, retained
witness information, or construction, explained below and in
\cref{app:lean-diff,app:lean-complexity}. The last nine rows concern
\cref{sec:complexity,app:small-arity}.}\label{tab:lean}\\
\toprule
Paper result & Formal counterpart\\\midrule\endfirsthead
\toprule Paper result & Formal counterpart\\\midrule\endhead
\bottomrule\endfoot
\eqref{eq:g3}--\eqref{eq:g4}, domains & \lean{Characterization.treeLe_eq_gfp}, \lean{Characterization.unfold_range_iff}, \lean{Characterization.toTree_range_iff}\\
Normalization (\cref{sec:normal}) & \lean{treeLe_iff_normalize}\\
Unsafe words, \eqref{eq:g13} & \lean{entails_iff_not_sideUnsafe}\\
Signed solutions (\cref{lem:signed}) & \lean{fixed_solution_correspondence}, \lean{signedDual_involutive}, \lean{SignSelection.sat_signedDual_of_flipClosed}, \lean{normalize_injective}\\
Selection (\cref{lem:selector}) & \lean{SignSelection.leastShape_sameShape_dual}, \lean{SignSelection.select_sat_of_coherent}, \lean{SignSelection.select_fixed}, \lean{SignSelection.select_leastShape_sat_fixed}\\
Signed satisfiability, \eqref{eq:g19} & \lean{sat_iff_signed_sat}\\
Fixed witnesses (\cref{lem:fixed-witness,thm:four-spine}) & \lean{Events.sideUnsafe_iff_fixedWitness}, \lean{Spine.Closure.fixedWitness_iff_not_labelClash}, \lean{Events.sideUnsafe_iff_not_labelClash}\\
Closure and path rules (\cref{sec:closure,lem:path-rules}) & \lean{Derives.sound}, \lean{UpperAt.comp}, \lean{LowerAt.comp}, \lean{UpperAt.le_embed}, \lean{LowerAt.embed_le}, \lean{UpperAt.factor}, \lean{LowerAt.factor}, \lean{LowerAt.decompose_upper}, \lean{LowerAt.decompose_lower}\\
Sign purity (\cref{lem:sign-purity}) & \lean{Signs.derives_sign}, \lean{Signs.upperAt_sign}, \lean{Signs.lowerAt_sign}\\
Satisfiability (\cref{thm:sat}) & \lean{leastShape_sat}, \lean{satisfiable_iff_not_labelClash}\\
Spine extension$^\dagger$ (\cref{lem:spine}) & \lean{Spine.Closure.witness_of_sat}, \lean{Spine.Closure.extend_sat} (both for fixed assignments), \lean{Spine.Closure.extension_flipClosed}, \lean{Spine.Closure.extension_signCoherent}\\
Separation and bridges (\cref{lem:separation,lem:bridge}) & \lean{Ranked.Extension.derives_into_source}, \lean{Ranked.Extension.derives_out_of_sink}, \lean{Ranked.Extension.derives_original} via \lean{Spine.Closure.extension_ranked}, \lean{Spine.Closure.derives_inv}, \lean{Spine.Closure.derives_lv_of_uAt}, \lean{Spine.Closure.derives_up_of_lAt}, \lean{Spine.Closure.derives_bridge}\\
Clashes and events (\cref{thm:clashes}) & \lean{Spine.Closure.labelClash_extension_iff}, \lean{Events.sideUnsafe_iff_not_events}\\
Label states and run (\cref{sec:label-reader}) & \lean{Monitor.card_state}, \lean{Monitor.mem_eval_U}, \lean{Monitor.mem_eval_L}, \lean{Monitor.eval_b_iff}\\
Rejected-path property (\cref{sec:interface}) & \lean{Monitor.rejectedPath}\\
Unsafe words are holes (\cref{thm:holes}) & \lean{Monitor.hole_iff_not_occurs}, \lean{Monitor.hole_iff_unsafe}\\
Tail (\cref{thm:brt}) & \lean{RejectedTail.boundedRejectedTail_of_rejectedPath}\\
Return gap (\cref{lem:return-gap}) & \lean{RejectedTail.return_gap}\\
Alignment (\cref{lem:alignment}) & \lean{LetteredHierarchy.alignment}\\
Cores (\cref{lem:cores}) & \lean{CoreOccurrences.corePos_succ_iff_occurrence}, \lean{LetteredHierarchy.startCore_invariants}, \lean{RejectedTail.cores_succ_subset}, \lean{StartCores.sCount_le_level}, \lean{CoreOccurrences.rho_at_greatest}\\
Equal-support width (\cref{lem:width}) & \lean{RejectedTail.equal_support_width}\\
Depth (\cref{thm:depth}) & \lean{RejectedTail.rejected_depth}, \lean{RejectedTail.depth_le_of_support_width}\\
Derivatives (\cref{lem:der-holes}) & \lean{LetteredHierarchy.isHole_der_iff}, \lean{LetteredHierarchy.isHole_hierOf_iff}, \lean{BoundedDepth.card_cores}\\
Comparison maps (\cref{lem:comparison-map}) & \lean{CoAlignment.isComp_map_iff}, \lean{LevelMaps.NoCollapse.reflect_last}\\
Exchange types (\cref{lem:exchange-types}) & \lean{Exchange.exType}, \lean{Exchange.card_exType_le}, \lean{NormalForm.normalTypeBound}, \lean{Exchange.Pset}, \lean{Exchange.Uset}\\
Isolated exchange (\cref{lem:isolated-exchange}) & \lean{Exchange.Isolated}, \lean{Exchange.isolated_merge}, \lean{NormalForm.isolated_merge_spec}\\
Level maps and start counts (\cref{lem:level-maps}) & \lean{LevelMaps.levelMap_spec}, \lean{LevelMaps.start_below}, \lean{LevelMaps.start_cap_up}, \lean{LevelMaps.LevelMapResult.depth_eq}\\
Dictionary normal form (\cref{lem:dictionary-normal-form}) & \lean{NormalForm.Normal}, \lean{NormalForm.dictionaryPotential_lt}, \lean{NormalForm.exists_normal_form}\\
Dictionary counting (\cref{lem:dictionary-count}) & \lean{NormalForm.normalAt_alphCard_le}, \lean{NormalForm.normalAllowance}\\
Propagation (\eqref{eq:boundary-propagation}) & \lean{BoundaryPairs.boundaryPairs_expand_subset}, \lean{BoundaryPairs.card_endpointPairs_le}, \lean{BoundaryPairs.boundaryCount_step_add_one}\\
Dictionary bounds (\cref{lem:dictionary-bound}) & \lean{ElementaryBound.normal_boundary_count_le}, \lean{ElementaryBound.normal_alphCard_le}, \lean{ElementaryBound.exists_bounded_dictionary_hole}\\
Agreement (\eqref{eq:monitor-agreement}, \cref{lem:agreement-derivative}) & \lean{BoundedDepth.Agree}, \lean{BoundedDepth.Agree.comp}, \lean{BoundedDepth.Agree.isHole_iff}, \lean{BoundedDepth.Agree.derive}\\
Column summary, \eqref{eq:col} & \lean{BoundedDepth.colEval}, \lean{BoundedDepth.der_type_of_summ}, \lean{BoundedDepth.exists_short_letters}, \lean{BoundedDepth.colBound}, \lean{BoundedDepth.card_CS}\\
Supported compression (\cref{lem:simultaneous-compression}) & \lean{BoundedDepth.PositiveDictionaryBound}, \lean{BoundedDepth.mem_lower_dictionary}, \lean{BoundedDepth.shorten_level_supported}, \lean{BoundedDepth.descentFactor}\\
Compression (\cref{thm:search}) & \lean{BoundedDepth.exists_short_hole_of_dictionary}, \lean{ElementaryBound.exists_short_hole}, \lean{ElementaryBound.lengthBound}\\
Closed form (\cref{lem:elementary-estimate}) & \lean{ElementaryBound.dictionaryBound_le}, \lean{ElementaryBound.lengthBound_le}, \lean{RejectedTail.holeLengthBound_le}\\
Finite search (\cref{thm:hole-decision}) & \lean{RejectedTail.hole_iff_bounded_length}, \lean{RejectedTail.decideHole}\\
Regular transfer$^\dagger$ (\cref{sec:regular}) & \lean{unfold_shapeGraph}, \lean{regular_fixed_solution}, \lean{regular_countermodel}, \lean{RegularTransfer.countermodel_spec}, \lean{entailsReg_iff_entails}\\
Finite satisfiability (\cref{thm:finite-sat,lem:finite-signed}) & \lean{satFin_iff}, \lean{satFin_iff_signed}, \lean{satFinB_iff}\\
Finite transfer$^\dagger$ (\cref{thm:finite-transfer}) & \lean{Ranked.Extension.cycleClash_iff}, \lean{FiniteTransfer.finite_fixedWitness}, \lean{FiniteTransfer.finite_countermodel_of_unsafe}, \lean{finite_countermodel_transfer}, \lean{entailsFin_iff}\\
Decision (\cref{thm:correct,thm:main}) & \lean{decideEntails}, \lean{decideEntailsReg}, \lean{decideEntailsFin}, and their correctness theorems\\\midrule
Summary monitor (\cref{lem:summary-monitor}) & \lean{Complexity.Summary.derivedTable}, \lean{Complexity.Summary.derivedTable_agree}\\
Summary search$^\dagger$ (\cref{prop:summary-search}) & \lean{Complexity.Summary.Config}, \lean{Complexity.Summary.summarySearch}, \lean{Complexity.Summary.summarySearch_iff_reachable}, \lean{Complexity.Summary.summarySearch_sound}, \lean{Complexity.Summary.summarySearch_complete}, \lean{Complexity.Summary.card_config_le}, \lean{Complexity.Summary.configBits_le}\\
Upper bounds$^\dagger$ (\cref{thm:complexity}) & \lean{Complexity.Summary.hole_iff_summarySearch}, \lean{Complexity.Summary.card_config_label_le}, \lean{decideEntailsSummary}, \lean{decideEntailsRegSummary}, \lean{decideEntailsFinSummary}, and their correctness theorems\\
One-letter readers$^\dagger$ (\cref{lem:one-letter}) & \lean{Complexity.Unary.isReaderHole_take}, \lean{Complexity.Unary.firstRejected_lt_card}, \lean{Complexity.Unary.exists_hole_length_le}, \lean{Complexity.Unary.hole_iff_length_le_card_sub_one}\\
Arity one (\cref{thm:arity-one}) & \lean{Complexity.Unary.unsafe_iff_bounded}, \lean{Complexity.Unary.firstRejectedCut}, \lean{Complexity.Unary.hole_iff_firstRejectedCut}, \lean{decideEntailsUnary}, \lean{decideEntailsRegUnary}, \lean{decideEntailsFinUnary}, and their correctness theorems\\
Top-free trees, item (3), and item (4) for $a\in\{0,1\}$ (\cref{lem:cx-topfree}) & \lean{Complexity.Rejection.trace_confined}, \lean{Complexity.FiniteTrees.trunc}, \lean{Complexity.FiniteTrees.node_letterChildren_le_trunc}\\
Counter automaton (\cref{lem:cx-counter}) & \lean{Complexity.Counter.exists_runs_iff}, \lean{Complexity.Counter.counterRecord}, \lean{Complexity.Counter.length_counterRecord}, \lean{Complexity.Counter.counterNFA}, \lean{Complexity.Counter.counterNFA_final}, \lean{Complexity.Counter.counterNFA_accepts_iff}, \lean{Complexity.Counter.length_counterWord}\\
Exponentially long holes$^\dagger$ (\cref{prop:cx-length}) & \lean{Complexity.rejectionSystem}, \lean{Complexity.Rejection.not_unsafe_left}, \lean{Complexity.Rejection.unsafe_iff}, \lean{Complexity.expSystem}, \lean{Complexity.satisfiableFin_expSystem}, \lean{Complexity.not_entails_expSystem}, \lean{Complexity.not_entailsReg_expSystem}, \lean{Complexity.not_entailsFin_expSystem}, \lean{Complexity.expSystem_unsafe_iff}, \lean{Complexity.expSystem_hole_iff}, \lean{Complexity.expSystem_shortest_hole}, \lean{Complexity.expSystem_shortest_refutesB}, \lean{Complexity.two_pow_le_length_counterWord}\\
PSPACE-hardness$^\dagger$ (\cref{thm:cx-hard}) & \lean{Complexity.Hardness.doublingCode}, \lean{Complexity.Hardness.eq_of_doublingCode_prefix}, \lean{Complexity.hardnessSystem}, \lean{Complexity.entails_hardnessSystem_iff}, \lean{Complexity.entailsReg_hardnessSystem_iff}, \lean{Complexity.entailsFin_hardnessSystem_iff}, \lean{Complexity.Hardness.numVars_eq}, \lean{Complexity.hardnessSystem_length_le}, \lean{Complexity.satisfiable_hardnessSystem}\\
\end{longtable}

The minimal-constructor-position result (\cref{lem:minimal-shape}), the
equivalence and alphabet reduction in \cref{sec:general-holes}, and the
countermodel size bounds in \cref{sec:countermodel-bounds} are not
mechanized. The explicit elementary estimate in \eqref{eq:elementary-closed}
is mechanized. The parts of \cref{sec:complexity} that are not mechanized
are listed in \cref{app:lean-complexity}.

\lean{RegularTransfer.countermodel_spec} records failure at the same word
as a disjunction of the two sides. The regular satisfiability equivalence
at the end of \cref{sec:regular} follows from the paper's construction;
there is no separately exported declaration of that equivalence.
For finite trees, \lean{FiniteTransfer.finite_fixedWitness} retains the
word in the signed witness, whereas the decoded export
\lean{FiniteTransfer.finite_countermodel_of_unsafe} records only order failure.

\subsection{Paper and Formalization}
\label{app:lean-diff}
Both presentations work directly over the child alphabet and use the same
label-state count and final rejected-tail, depth, and search bounds. All three
domains use the same four-spine construction and selected least shape.
The paper and formalization differ in representation, in the strength
of some intermediate statements, and in intermediate numerical bounds.
\begin{itemize}[leftmargin=*]
\item \emph{Representation.} The paper uses $V\times\{0,1\}$ for signed variables and freely names
fresh spine variables. The code builds $\Psi^4_w$ directly in the signed
finite ordinal type of size $2(k+2|w|+2)$ and generates its upper half by
flipping the lower half. Its old-variable map is an injective embedding.
The general \lean{Ranked.Extension} accepts this embedding, rather than
requiring the old signed variables to occupy an initial segment;
\lean{Spine.Closure.extension_ranked} supplies the instance. Its rank is
the base-variable index, whereas the paper uses $0$, $j$, and $|w|+1$
for old variables, spine cuts, and fillers.
\item \emph{Proof organization.} The paper packages the bridge formula as an equivalence.
The code uses an invariant for forward closure derivations and separate
backward construction lemmas; the correspondence table lists both.
\item \emph{Intermediate statement strength.} The paper proves both directions of
\cref{lem:spine} for any old assignment satisfying the base system.
The code's extension and restriction interface is
stated for fixed assignments, which suffices for every witness transfer.
The paper assigns suffix embeddings to fresh spine variables;
\lean{Spine.Closure.extend} and \lean{Spine.Closure.extendBase} instead
extend the fixed witness using traces of its root trees.
\item \emph{Hypotheses.} The clash equivalence
\lean{Spine.Closure.labelClash_extension_iff} holds for any flip-closed,
sign-coherent system and any roots $X,Y$. The fixed-witness equivalence
\lean{Spine.Closure.fixedWitness_iff_not_labelClash} has the same scope.
The formal reader and rejected-path property require neither flip closure
nor sign coherence.
The inhabited-alphabet assumptions of intermediate depth and compression
results are removed by the final bounded-length theorem, as explained
in \cref{app:lean-decision}.
\item \emph{Graph representation.} For the least-shape graph, the paper stores upper and lower sets.
The code computes the upper set as a lower set in the order-dual system;
these are the same supporting-variable sets by path duality.
Decoding uses \lean{decodedGraph}, whose normalization
\lean{RGraph.normalize} of a graph $g$ is \lean{RGraph.mix} applied to $g$
twice, with $2|g|^2$ states and only diagonal graph pairs reachable, whereas the paper uses a single graph state and a polarity bit;
the size and depth bounds are not exported theorems.
\item \emph{Elementary compression.} The functions $C,A,G$ and $\holeBound$
match \lean{NormalForm.normalTypeBound}, \lean{NormalForm.normalAllowance},
\lean{ElementaryBound.dictionaryBound}, and
\lean{ElementaryBound.lengthBound}. The paper uses finite sets for
recursive supports; \lean{BoundedDepth.shorten_level_supported} uses
lists with a no-duplicates proof and explicitly calls \texttt{dedup}.
The fixed-original-dictionary invariant is the same. The paper's ordered
boundary-pair inequality subtracts one; the formal recurrence uses
\lean{BoundaryPairs.boundaryCount_step_add_one}, with one added to the
left side, avoiding natural-number subtraction. The normalization maps
preserve geometry and start counts; neither presentation assumes global
agreement of higher derived predicates after a merge. The stronger
all-state agreement is used separately for column compression.
\item \emph{Periodicity proofs in the tail bound.} The paper's branching case uses a rotation
argument where the code uses Fine--Wilf periodicity \citep{FineWilf65}. Both prove the same
intermediate bounds, $4N-2$ for a branching cycle and $2N$ for a forced cycle.
\end{itemize}


The separate kernel-checked formalization of the NP03 counterexample
agrees with the argument in \cref{app:np03}. It is not part of this
artifact or a dependency of the decision procedures.

\subsection{Complexity Results}
\label{app:lean-complexity}
The module \texttt{Complexity} formalizes the combinatorial content of
\cref{sec:complexity} and of \cref{app:small-arity}; the last nine rows of
\cref{tab:lean} give the correspondence. Importing \texttt{DeciNSSE} also
exposes six further procedures. \lean{decideEntailsSummary}, \lean{decideEntailsRegSummary}, and
\lean{decideEntailsFinSummary} decide the three entailment relations for
every arity and variance by the summary search, at the depth bound
\lean{RejectedTail.holeDepthBound} and the dictionary bound
\lean{ElementaryBound.dictionaryBound} at that depth; at arity zero their
hole tests examine only the empty word
(\lean{Complexity.Summary.hole_iff_nil_of_eq_zero}).
\lean{decideEntailsUnary}, \lean{decideEntailsRegUnary}, and
\lean{decideEntailsFinUnary} decide them at arity one, by the test of
\cref{thm:arity-one}. Each has a correctness theorem of the form of
\lean{decideEntails_correct}, without premises. Beyond the trusted statement
of \cref{app:lean-statement}, the hardness theorems and the theorems on
exponentially long holes depend on the definitions of binary automata and
their languages (\lean{Complexity.BinNFA},
\lean{Complexity.BinNFA.Accepts}) and of the constructed systems.

The complexity classes themselves are not mechanized: space-bounded
machines, the space accounting of \cref{app:complexity}, Savitch's theorem,
and the space hierarchy theorem. Nor are the PSPACE-hardness of
universality, which \cref{thm:cx-hard} cites, the polynomial-time tests of
\cref{lem:poly-tests}, the polynomial-time computability of the reduction and
of $\mathcal B_\beta$, the polynomial-time procedure of
\cref{prop:arity-zero} with its general systems, \cref{prop:cx-embedding},
and \cref{ex:cx-padding}. Items (1) and (2) of \cref{lem:cx-topfree} are
not stated separately. The formalization differs from the paper as follows.
\begin{itemize}[leftmargin=*]
\item \emph{Summary search.} \Cref{prop:summary-search} describes a
nondeterministic search that counts the columns of the guessed column word
and stops before $\mathit{col}(N+j,h)$ of them.
\lean{Complexity.Summary.summarySearch} instead decides reachability of an
accepting configuration in a finite graph
(\lean{Complexity.Summary.summarySearch_iff_reachable}), which it computes
by saturating the set of reachable configurations
(\lean{Complexity.reachSet}). A configuration,
\lean{Complexity.Summary.Config}, holds the level, the number of current
letters, the current table, and, during a derive move, the marker, the
number of blocks, the summary, and the track flags. It has no column counter, since reachability in a finite
graph needs no bound on the length of a path. Hence completeness uses
neither column shortening nor \cref{lem:summary-agreement}: the search reads
the whole column word of the lifted blocks, and
\lean{Complexity.Summary.derivedTable_agree} composed with
\lean{BoundedDepth.Agree.derive} gives the agreement at the next level. For a
monitor with $N$ states over $n$ letters, at most $g$ blocks per level, and
depth bound $d$, \lean{Complexity.Summary.card_config_le} and
\lean{Complexity.Summary.configBits_le} bound the number of configurations
by $2^{11(N+d+1)^2(n+g+1)^2}$, matching the space bound of
\cref{prop:summary-search} with $|\Gamma|=n$. For the label readers,
\lean{Complexity.Summary.card_config_label_le} gives the bound
$2^{(n+1)^2\cdot2^{2^{20k+20}}}$, so a configuration fits in
$(n+1)^2\cdot2^{2^{20k+20}}$ bits, that is, in
$2^{2^{O(k)}}\cdot\operatorname{poly}(n)$ bits, as in the proof of
\cref{thm:complexity}. These bounds support the space argument; the
saturation itself is not proved to run within that space.
\item \emph{Hardness.} The paper treats an automaton with one initial state
$q_0$, sets $x:=X_{q_0}$, and builds $\chi(\mathcal A')$ from the code
automaton $\mathcal A'$. \lean{Complexity.hardnessSystem} is built from the
automaton directly and allows several initial states. Its transition
variables $Y_t$ stand for the middle states of $\mathcal A'$, and $b$ stands
for $X_{q_{\mathrm f}}$, which is $\Bot$ in every solution. Flat literals have
no inequality between two variables, so a sandwich variable $z$, with the
literals $f(x,b,\ldots,b)\le z$ and $z\le f(X_p,b,\ldots,b)$ for every
initial state $p$, expresses that $x\le X_p$ for all of them: by
decomposition, such a $z$ exists iff these inequalities hold. Both proofs
therefore establish the same theorem, for every $n\ge2$ and each domain: an
always satisfiable covariant system of linear size entails $x\le r$ iff the
automaton is universal. With the PSPACE-hardness of universality, either
gives \cref{thm:cx-hard}. The universal
case of the paper's proof
uses the identification $x=X_{q_0}$ only through $\sigma(x)\le\sigma(X_{q_0})$,
and in the rejecting case $z=f(\sigma(x),\Bot,\ldots,\Bot)$ completes the
countermodel. With several initial states, the reduction to one initial
state in \cref{sec:complexity-lower} is not needed. For an automaton with
the states $Q$ and a list $\delta$ of transitions, the formal system has
$\len Q+\len\delta+6$ variables and at most $2\len Q+2\len\delta+4$ literals
(\lean{Complexity.Hardness.numVars_eq},
\lean{Complexity.hardnessSystem_length_le}); the paper's has
$\len{Q'}+4=\len Q+\len\delta+5$ variables.
\item \emph{Exponentially long holes.} \lean{Complexity.rejectionSystem} produces the
literals of $\Phi_\beta^{(n)}$, up to repetition, from any automaton all of
whose states are final. Its variables $x$, $g$, and $b$ are named as in the
paper, and its variable $y$ receives the lower literals of the initial
states directly, as $Y_{q_{\mathrm{in}}}$ does. The semantic theorems
\lean{Complexity.Rejection.unsafe_iff} and \lean{Complexity.Rejection.hole_iff}
hold for every such automaton whose language is the set of words without a
fixed nonempty prefix, and \lean{Complexity.expSystem} applies them to
\lean{Complexity.Counter.counterNFA}. In \lean{Complexity.expSystem}, the
Lean parameter named \texttt{k} corresponds to $\beta$, and the variable count
$9\beta+31$ is the paper's $k$. The literal count is only bounded, by $81\beta+247$
(\lean{Complexity.expSystem_length_le}).
\item \emph{One-letter readers.} \lean{Complexity.Unary.exists_hole_length_le}
holds for every reader over an alphabet with at most one letter, which
includes the empty alphabet, whereas \cref{lem:one-letter} assumes one
letter; neither needs the rejected-path property. As in the proof of
\cref{thm:arity-one}, the unary procedures follow the run to the first
rejected cut $J<N$ (\lean{Complexity.Unary.firstRejectedCut}) and test the
single word $0^J$ (\lean{Complexity.Unary.hole_iff_firstRejectedCut}).
\item \emph{Arity $n\ge3$.} The formalization proves the hardness theorems
and the theorems on exponentially long holes directly for every $n\ge2$, with $b$ in the arguments
$2,\ldots,n-1$, and does not use \cref{prop:cx-embedding}. The paper proves
\cref{thm:cx-hard} directly as well, but reaches $n\ge3$ in
\cref{prop:cx-length} through the embedding.
\end{itemize}
